\documentclass{article}
\usepackage{anyfontsize}
\usepackage[UTF8]{ctex} % 如果需要支持中文
\usepackage{fontspec} 

\usepackage[a4paper, left=0.1in, right=0.1in, top=0.05in, bottom=0.05in]{geometry} % 设置四边页边距为0.1英寸{geometry} % 设置页边距为0.5英寸
\usepackage{enumitem} % 用于自定义列表
\usepackage{amsmath}  % 数学公式支持

\setlength{\parindent}{0pt} % 不缩进
\usepackage{etoolbox} % 用于列表操作
%调整类代码
\usepackage{comment}
\usepackage{tocloft} % 用于定制目录样式
\usepackage{hyperref} %不需要目录盒子注释这
\usepackage{ifthen}
\usepackage{tikz}
\usepackage{titletoc}
\usepackage{graphicx}
\usepackage{amssymb}  % 提供数学符号，包括\therefore
%目录设定
%快捷键 CMD + / 注释
%  \renewcommand{\cftdot}{} % 隐藏目录中的页码
%  \cftpagenumbersoff{section}
%  \cftpagenumbersoff{subsection}
%  \makeatletter
%  \renewcommand{\@pnumwidth}{0em}  % 设置页码宽度为0
%  \renewcommand{\@tocrmarg}{0em}   % 设置右边距为0
%  \makeatother
 
%\setCJKmainfont[
 %   Path = C:/USERS/11252/APPDATA/LOCAL/MICROSOFT/WINDOWS/FONTS/,
 %   UprightFont = {SourceHanSerifCN-Light-5.otf},
 %   BoldFont = {SourceHanSerifCN-Medium-6.otf},
%    Scale=1
%]{SourceHanSerifCN}

%\setmainfont{SourceHanSerif}  % 设置英文字体

% 处理冲突的命令
\let\oldbackepsilon\backepsilon
\let\backepsilon\relax
\let\oldeth\eth
\let\eth\relax
\let\olddigamma\digamma
\let\digamma\relax

\usepackage{float}
\usepackage[a4paper, left=0.1in, right=0.1in, top=0.05in, bottom=0.05in]{geometry} % 设置四边页边距为0.1英寸{geometry} % 设置页边距为0.5英寸
\usepackage{enumitem} % 用于自定义列表
\usepackage{amsmath}  % 数学公式支持

\setlength{\parindent}{0pt} % 不缩进
\usepackage{etoolbox} % 用于列表操作
%调整类代码
\usepackage{comment}
\usepackage{tocloft} % 用于定制目录 样式
\usepackage{hyperref} %不需要目录盒子注释这
\usepackage{ifthen}
\usepackage{tikz}
\usepackage{titletoc}
\usepackage{graphicx}
\usepackage{amssymb}  % 提供数学符号，包括\therefore
%目录设定
%快捷键 CMD + / 注释
%  \renewcommand{\cftdot}{} % 隐藏目录中的页码
%  \cftpagenumbersoff{section}
%  \cftpagenumbersoff{subsection}
%  \makeatletter
%  \renewcommand{\@pnumwidth}{0em}  % 设置页码宽度为0
%  \renewcommand{\@tocrmarg}{0em}   % 设置右边距为0
%  \makeatother
 


\usepackage{environ}

\newif\ifshowcontent  % 定义一个条件判断变量
\showcontenttrue    % 设置为 false，隐藏所有 question 环境的内容

\newcounter{question} % 题目计数器
\setcounter{question}{0}
\NewEnviron{question}{%
    \ifshowcontent
        \par\stepcounter{question}\textbf{\thequestion.} \BODY\par % 如果显示内容，则输出
    \fi
}

\NewEnviron{choices}{%
    \ifshowcontent
        \begin{enumerate}[label=\Alph*., leftmargin=*]
            \BODY
        \end{enumerate}\par
    \fi
}

\NewEnviron{correctanswer}{%
    \ifshowcontent
        \textbf{答案:}\BODY\par
    \fi
}



\newenvironment{explanation}
{\textbf{解析:}\par}
{\par}



% 新增考察方面命令
\newcommand{\checklist}[1]{
    \textbf{考察:} #1 \par % 在题目下方显示考察方面
    \addcontentsline{toc}{subsection}{题号\thequestion 考察: #1} % 将考察内容添加到目录
    \vspace{2em} % 添加1em的垂直空间
}

\graphicspath{{images/}} %统一

\begin{document}
 %\fontsize{22}{31}\selectfont
 \tableofcontents % 添加目录
\newpage
%\pagestyle{empty}




\section{考点1:信用风险基础}
\setcounter{question}{0}


\begin{question}
    Sarah Chen, FRM, is a bank credit analyst who is examining the financial statements of a bank. She notices that there is a paragraph noted in the auditor's report that states that although the auditors agreed with virtually all of the bank's accounting treatments of the financial statement items, the auditors did not agree with the bank's decision to treat some of the leases as operating leases instead of capital leases. Based on that information, which of the following audit report opinions has the auditor most likely issued?
\end{question}

\begin{choices}
    \item Adverse opinion
    \item Denial of opinion
    \item Qualified opinion
    \item Unqualified opinion
\end{choices}

\begin{correctanswer}
    C
\end{correctanswer}

\begin{explanation}
    \textbf{A选项错误。}
    \begin{itemize}
        \item 否定意见（Adverse opinion）表示财务报表存在全局性、系统性错误的情况，而本题中只有一个具体的会计处理问题。
    \end{itemize}
    
    \textbf{B选项错误。}
    \begin{itemize}
        \item 拒绝表示意见（Denial of opinion）表示审计师无法获取足够的审计证据，而本题中审计师已经完成了审计工作。
    \end{itemize}
    
    \textbf{C选项正确。}
    \begin{itemize}
        \item 有保留意见（Qualified opinion）适用于审计师对财务报表整体无异议，但对特定项目持不同意见。本题中审计师对租赁会计处理有不同意见，但认可其他会计处理，因此最可能出具保留意见。
    \end{itemize}
    
    \textbf{D选项错误。}
    \begin{itemize}
        \item 无保留意见（Unqualified opinion）表示审计师完全认可财务报表，而本题中审计师对租赁会计处理有不同意见。
    \end{itemize}
\end{explanation}

\checklist{审计意见的差异源于问题的“严重性”和“影响范围”（考概念辨析）}


\begin{question}
    Bank of the Midwest has been struggling with poor asset quality for some time. The bank lends primarily to large farming operations that have struggled in recent years due to a glut of wheat and cotton on the market. Bank regulators have recently required that the bank write off some of these loans, which has entirely wiped out the capital of the bank. However, the bank still has some liquidity sources it can use, including a correspondent bank and the Federal Reserve. Bank of the Midwest is:
\end{question}

\begin{choices}
    \item An insolvent but not failed bank
    \item Both a failed bank and an insolvent bank
    \item Neither a failed bank nor an insolvent bank
    \item A failed bank but not an insolvent bank
\end{choices}

\begin{correctanswer}
    A
\end{correctanswer}

\begin{explanation}
    \textbf{A选项正确。}
    \begin{itemize}
        \item 题目明确说明银行的资本被完全冲销（Insolvent），但仍有流动性来源维持运营 → 未达到破产条件。
        \item 典型的“资不抵债但未破产”场景
    \end{itemize}
    
    \textbf{B选项错误。}
    \begin{itemize}
        \item 该银行虽然资不抵债，但由于仍能获得流动性继续运营，所以并未破产。
    \end{itemize}
    
    \textbf{C选项错误。}
    \begin{itemize}
        \item 该银行确实资不抵债，因为其资本已被完全冲销。
    \end{itemize}
    
    \textbf{D选项错误。}
    \begin{itemize}
        \item 逻辑矛盾。破产必然以资不抵债为前提。
    \end{itemize}
\end{explanation}

\checklist{资不抵债（资本耗尽） vs. 破产（丧失持续经营能力）细节辨析题} 


\begin{question}
    In regard to bank failures, each of the following is correct, except which is false?
\end{question}

\begin{choices}
    \item If liabilities exceed assets, a firm is insolvent
    \item Although bank failures are very common, bank insolvencies are historically very rare
    \item Bankruptcy is a legal status, and is not implied by (is not a synonym for) insolvency
    \item It is possible for an insolvent bank to continue to operate
\end{choices}

\begin{correctanswer}
    B
\end{correctanswer}

\begin{explanation}
    \textbf{A选项正确，不选。}
    \begin{itemize}
        \item 资不抵债（Insolvency）：负债 > 资产
    \end{itemize}
    
    \textbf{B选项正确。}
    \begin{itemize}
        \item 实际情况相反：银行在资不抵债后可能继续运营（如获得流动性支持），因此资不抵债比破产更常见。
        \item 监管机构通常在资不抵债阶段介入，以避免破产，所以银行破产的数量通常少于资不抵债的银行。要仍能获得流动性支持，就可以继续运营。
    \end{itemize}
    
    \textbf{C选项正确，不选。}
    \begin{itemize}
        \item 教材强调资不抵债是财务状态，而破产是法律程序，两者不同。
    \end{itemize}
    
    \textbf{D选项正确，不选。}
    \begin{itemize}
        \item 银行在资不抵债后仍可通过外部流动性（如央行支持）维持运营。
    \end{itemize}
\end{explanation}

\checklist{区分资不抵债和破产（财务状态 vs. 法律程序）}


\begin{question}
    David Wang, a credit analyst with City Bank, is considering the loan application of a small, local car dealership. The dealership has been solely owned by John Smith for more than 20 years and sells three brands of American automobiles. Because of the suburban location, most of the cars sold in the past by the dealership have been large pick-up trucks and sports utility vehicles. However, sales have declined, and gasoline prices have continued to increase. As a result, Smith is considering selling a line of electric cars. Smith has borrowed from City Bank before but currently does not have a balance outstanding with the bank. Which of the following statements is not one of the four components of credit analysis Wang should be evaluating when performing the credit analysis for this potential loan?
\end{question}

\begin{choices}
    \item The business environment, competition, and economic climate in the region
    \item Smith's character and past payment history with the bank
    \item The car dealership's balance sheets and income statements for the last few years as well as Smith's personal financial situation
    \item The financial health of Smith's friends and family who could be called upon to guarantee the loan
\end{choices}

\begin{correctanswer}
    D
\end{correctanswer}

\begin{explanation}
    \textbf{A选项错误。}
    \begin{itemize}
        \item 业务环境、竞争格局和区域经济气候直接影响借款人的还款能力，是是“信用分析四要素”中的Conditions（环境/条件）
    \end{itemize}
    
    \textbf{B选项错误。}
    \begin{itemize}
        \item 借款人的个人信用记录（如历史还款表现）和品性是“信用分析四要素”中的Character（品格）
    \end{itemize}
    
    \textbf{C选项错误。}
    \begin{itemize}
        \item 企业财务报表（资产负债表、利润表）和借款人个人财务状况是量化偿债能力的基础数据，是“信用分析四要素”中的Capacity（偿债能力）。
    \end{itemize}
    
    \textbf{D选项正确。}
    \begin{itemize}
        \item 第三方担保人（如亲友）的财务健康并非标准信用分析组件。
        \item 信贷决策应基于借款人自身资质，即使要求担保，担保人分析也属于附加程序，而非核心四要素
    \end{itemize}
\end{explanation}

\checklist{信用分析的核心四要素}


\begin{question}
    Michael Chen, FRM, is a rating agency analyst who is currently performing financial statement analysis on a major bank. Which of the following financial statements would be least useful for bank credit analysis?
\end{question}

\begin{choices}
    \item Balance sheet
    \item Income statement
    \item Statement of cash flows
    \item Statement of changes in capital funds
\end{choices}

\begin{correctanswer}
    C
\end{correctanswer}

\begin{explanation}
    \textbf{A选项错误。}
    \begin{itemize}
        \item 资产负债表是最基础最重要的财务报表，是评估银行偿付能力的基础。
        \item 能显示银行的资本充足率、贷款资产、存款负债结构、流动性风险等。
    \end{itemize}
    
    \textbf{B选项错误。}
    \begin{itemize}
        \item 利润表是银行评级的关键维度。
        \item 反映银行的盈利能力、收入构成和成本结构。
    \end{itemize}
    
    \textbf{C选项正确。}
    \begin{itemize}
        \item 现金流量表对银行信用分析的用处最小。
        \item 现金流量表在分析非金融实体时很有用（可以区分经营、投资和筹资活动的现金流入和流出）
        \item 但银行业务本质上是资金中介，现金流管理方式特殊，所以现金流量表信息相对不那么重要。
    \end{itemize}
    
    \textbf{D选项错误。}
    \begin{itemize}
        \item 资本变动表是银行抗风险能力的重要指标。
        \item 反映资本充足率变化、股息政策和资本补充能力。
    \end{itemize}
\end{explanation}

\checklist{如何使用银行信用分析时各财务报表的作用}




\newpage



\section{考点2:信用风险管治}
\setcounter{question}{0}


\begin{question}
    In regard to bank failures, each of the following is correct, except which is false?
\end{question}

\begin{choices}
    \item If liabilities exceed assets, a firm is insolvent
    \item Although bank failures are very common, bank insolvencies are historically very rare
    \item Bankruptcy is a legal status, and is not implied by (is not a synonym for) insolvency
    \item It is possible for an insolvent bank to continue to operate
\end{choices}

\begin{correctanswer}
    B
\end{correctanswer}

\begin{explanation}
    \textbf{A选项正确，不选。}
    \begin{itemize}
        \item 资不抵债的经典定义就是"负债超过资产"
    \end{itemize}
    
    \textbf{B选项错误。}
    \begin{itemize}
        \item 这个说法是错误的。实际上，银行破产（failures）相对较少，而银行资不抵债（insolvencies）则更为常见。
        \item 银行即使资不抵债，只要仍能获得流动性支持，就可以继续运营。
        \item 偿付能力衡量的是长期情况，而流动性看的是短期情况。
    \end{itemize}
    
    \textbf{C选项正确，不选。}
    \begin{itemize}
        \item 资不抵债是财务状态，破产是法律程序
        \item 资不抵债是破产的必要条件，但不是充分条件。
    \end{itemize}
    
    \textbf{D选项正确，不选。}
    \begin{itemize}
        \item 资不抵债的银行如果能够获得流动性支持（如来自其他银行或央行的资金），确实可以继续运营。
    \end{itemize}
\end{explanation}

\checklist{破产与资不抵债的区别、银行持续经营的监管特例（注意流动支持）}

\begin{question}
    Dragon Capital, Inc., made a loan to a small start-up firm. The firm grew rapidly, and it appeared that Dragon had made a good credit decision. However, the firm grew too fast and could not sustain the growth. It eventually failed. Dragon had initially estimated its exposure at default to be \$1,500,000. Because of the firm's rapid growth and resulting increases in the line of credit, Dragon ultimately lost \$1,800,000. In terms of credit risk, this is an example of:
\end{question}

\begin{choices}
    \item Default on payment for goods or services already rendered
    \item A more severe loss than expected due to a ratings downgrade by a rating agency
    \item A more severe loss than expected due to a greater than expected exposure at the time of a default
    \item A more severe loss than expected due to a lower than expected recovery at the time of a default
\end{choices}

\begin{correctanswer}
    C
\end{correctanswer}

\begin{explanation}
    \textbf{A选项错误。}
    \begin{itemize}
        \item 描述的是已交付商品或服务的付款违约，与本题中涉及的金融债务无关。
    \end{itemize}
    
    \textbf{B选项错误。}
    \begin{itemize}
        \item 涉及评级下调，但题干未提及评级变化
    \end{itemize}
    
    \textbf{C选项正确。}
    \begin{itemize}
        \item 信用风险中的非预期损失（UL）可能因EAD超预估而扩大。
        \item 本题中，因借款人信用额度提升（即EAD动态增加），最终损失超过初始估值，属于“违约时敞口大于预期”的情景。
    \end{itemize}
    
    \textbf{D选项错误。}
    \begin{itemize}
        \item 讨论回收率（Recovery Rate）低于预期，但题干未提及。
    \end{itemize}
\end{explanation}

\checklist{对风险敞口动态性与信用损失关系的理解}

\begin{question}
    A bank has booked a loan with total commitment of USD 60,000 of which 75\% is currently outstanding. The default probability of the loan is assumed to be 3\% for the next year and loss given default (LGD) is estimated at 40\%. The standard deviation of LGD is 35\% and the standard deviation of the default event indicator is 8\%. Drawdown on default is assumed to be 70\%. The expected losses for the bank are:
\end{question}

\begin{choices}
    \item USD 450
    \item USD 540
    \item USD 666
    \item USD 720
\end{choices}

\begin{correctanswer}
    C
\end{correctanswer}

\begin{explanation}
    \textbf{预期损失（Expected Loss, EL）计算公式：} 
    \[
    \text{EL} = \text{EAD} \times \text{PD} \times \text{LGD}
    \]
    
    \textbf{步骤一：计算违约敞口（EAD）}
    \begin{itemize}
        \item 总授信额度（Total Commitment）：\$60{,}000
        \item 已提款部分（Outstanding）：\$60{,}000 × 75\% = \$45{,}000
        \item 未提款部分中，违约时将提取 70\%：\$60{,}000 × 25\% × 70\% = \$10{,}500
        \item EAD = \$45{,}000 + \$10{,}500 = \$55{,}500
    \end{itemize}
    
    \textbf{步骤二：代入公式计算}
    \begin{align*}
    \text{EL} &= 55{,}500 \times 0.03 \times 0.4 \\
             &= 666
    \end{align*}
    
    
    \end{explanation}

\checklist{预期损失的计算公式（理解公式中的风险敞口，排除干扰数据）}


\begin{question}
    You have been asked by the Chief Risk Officer of your bank to determine how much should be set aside as a loan-loss reserve for a 1-year horizon on a USD 150 million line of credit that has been extended to a large corporate borrower. Of the original balance, USD 30 million has already been drawn and due to deteriorating economic conditions the bank is concerned that the borrower might find itself in a liquidity crisis causing it to draw on the remaining commitment and default. Given the following information from the bank's internal credit risk models, what is an appropriate loan loss reserve to cover this eventually? 1-year default probability = 0.45\% Drawdown given default = 85\% Loss given default = 55\%
\end{question}

\begin{choices}
    \item USD 280,000
    \item USD 327,000
    \item USD 224,000
    \item USD 196,000
\end{choices}

\begin{correctanswer}
    B
\end{correctanswer}
\begin{explanation}
    \textbf{预期损失（Expected Loss, EL）计算公式：} 
    \[
    \text{EL} = \text{EAD} \times \text{PD} \times \text{LGD}
    \]

    \textbf{已知条件：}
    \begin{itemize}
        \item 总授信额度：\$150,000,000
        \item 已提款：\$30,000,000
        \item 未提款：\$120,000,000
        \item 违约时提款比例（DGD）：85\%
        \item 违约概率（PD）：0.45\% = 0.0045
        \item 违约损失率（LGD）：55\% = 0.55
    \end{itemize}

    \textbf{计算过程：}
    \begin{itemize}
        \item 违约时提款金额 = \$120,000,000 × 85\% = \$102,000,000
        \item 风险敞口（EAD） = \$30,000,000 + \$102,000,000 = \$132,000,000
        \item 预期损失 = \$132,000,000 × 0.0045 × 0.55 = \$327,000
    \end{itemize}

    \textbf{正确答案：B}
\end{explanation}



\checklist{预期损失的计算公式（注意公式中风险敞口）}

\begin{question}
    The senior management of a large credit card issuer wants to strengthen the policies and procedures that the firm uses to manage its model risk. The CRO explains to the senior management the roles and responsibilities of different business functions in managing model risk in alignment with best practices and existing supervisory expectations. Which of the following statements would the CRO be correct to make?
\end{question}

\begin{choices}
    \item The enterprise risk management function should be responsible for developing and maintaining documentation for the firm's models
    \item The model development team should be responsible for maintaining the firm-wide framework for managing model risk
    \item The validation team should start the model validation process after a model has been put in use and has begun to generate observable data
    \item The internal audit function should verify that model owners are complying with the firm's policies for managing model risk
\end{choices}

\begin{correctanswer}
    D
\end{correctanswer}

\begin{explanation}
    \textbf{A选项错误。}
    \begin{itemize}
        \item 模型文档的开发和维护应由模型开发团队主导，而非风险管理部门。
        \item 风险管理（ERM）的职责是监督框架执行。
    \end{itemize}
    
    \textbf{B选项错误。}
    \begin{itemize}
        \item 全行模型风险管理框架的制定和维护是风险管理职能的核心职责，而非模型开发团队。
    \end{itemize}
    
    \textbf{C选项错误。}
    \begin{itemize}
        \item 监管要求验证必须前置，以确保模型投产前的可靠性。
        \item 所以模型验证应在投入使用前完成，而非依赖使用后生成的数据。
    \end{itemize}
    
    \textbf{D选项正确。}
    \begin{itemize}
        \item 内部审计作为第三道防线，负责核查模型负责人是否遵守模型风险管理相关制度。
    \end{itemize}
\end{explanation}

\checklist{模型风险管理的“三道防线”架构}





\newpage



\end{document}